% ============================================================
%  CUADERNILLO DE ESTUDIO — CLASE 3
%  Seguridad de la Información y de Sistemas
%  Lic. en Gestión de Negocios Digitales — UCA
%  Compilar con: pdflatex (dos pasadas para el índice)
% ============================================================

\documentclass[12pt,a4paper]{article}

% ---------- Paquetes ----------
\usepackage[utf8]{inputenc}
\usepackage[spanish]{babel}
\usepackage[T1]{fontenc}
\usepackage{lmodern}
\usepackage[margin=2.2cm,headheight=14pt]{geometry}
\usepackage{amsmath}
\usepackage{amssymb}
\usepackage{enumitem}
\usepackage{booktabs}
\usepackage{array}
\usepackage{tabularx}
\usepackage{xcolor}
\usepackage{tcolorbox}
\usepackage{graphicx}
\usepackage{fancyhdr}
\usepackage{titlesec}
\usepackage{hyperref}
\usepackage{multicol}
\usepackage{pifont}
\usepackage{wasysym}
\usepackage{tikz}
\usetikzlibrary{positioning,arrows.meta,shapes.geometric}

% ---------- Colores ----------
\definecolor{azulUCA}{HTML}{003366}
\definecolor{doradoUCA}{HTML}{B8860B}
\definecolor{grisClaro}{HTML}{F2F2F2}
\definecolor{rojoAlerta}{HTML}{C0392B}
\definecolor{verdeOK}{HTML}{27AE60}
\definecolor{azulCielo}{HTML}{2980B9}
\definecolor{naranjaAmenaza}{HTML}{D35400}

% ---------- tcolorbox estilos ----------
\tcbuselibrary{skins,breakable}

\newtcolorbox{cajaTitulo}[1][]{
  colback=azulUCA, colframe=azulUCA, coltext=white,
  fonttitle=\Large\bfseries, arc=4mm, #1
}

\newtcolorbox{cajaDefinicion}{
  colback=azulCielo!10, colframe=azulCielo,
  fonttitle=\bfseries, title=Definici\'on, arc=3mm, breakable
}

\newtcolorbox{cajaAmenaza}{
  colback=naranjaAmenaza!10, colframe=naranjaAmenaza,
  fonttitle=\bfseries\color{naranjaAmenaza},
  title={\ding{55}\; Amenaza}, arc=3mm, breakable
}

\newtcolorbox{cajaHerramienta}{
  colback=verdeOK!10, colframe=verdeOK,
  fonttitle=\bfseries\color{verdeOK},
  title={\ding{45}\; Herramienta}, arc=3mm, breakable
}

\newtcolorbox{cajaCaso}{
  colback=doradoUCA!10, colframe=doradoUCA,
  fonttitle=\bfseries, title=Caso de Estudio, arc=3mm, breakable
}

\newtcolorbox{cajaActividad}{
  colback=grisClaro, colframe=azulUCA,
  fonttitle=\bfseries\color{azulUCA},
  title={\ding{45}\; Actividad Pr\'actica}, arc=3mm, breakable
}

\newtcolorbox{cajaNota}{
  colback=rojoAlerta!8, colframe=rojoAlerta,
  fonttitle=\bfseries\color{rojoAlerta}, title=Nota Importante, arc=3mm, breakable
}

% ---------- Header / Footer ----------
\pagestyle{fancy}
\fancyhf{}
\fancyhead[L]{\small\textcolor{azulUCA}{Seguridad de la Informaci\'on y de Sistemas}}
\fancyhead[R]{\small\textcolor{azulUCA}{Cuadernillo --- Clase 3}}
\fancyfoot[C]{\thepage}
\fancyfoot[L]{\small\textcolor{gray}{UCA}}
\renewcommand{\headrulewidth}{0.4pt}
\renewcommand{\footrulewidth}{0.4pt}

% ---------- Títulos de sección ----------
\titleformat{\section}
  {\color{azulUCA}\Large\bfseries}
  {\thesection.}{0.6em}{}
\titleformat{\subsection}
  {\color{azulCielo}\large\bfseries}
  {\thesubsection}{0.6em}{}

% ---------- Enlaces ----------
\hypersetup{colorlinks=true, linkcolor=azulUCA, urlcolor=azulCielo}

% ---------- Checkbox ----------
\newcommand{\checkboxVacio}{$\square$\;}

% ============================================================
\begin{document}

% ============================================================
%  PORTADA
% ============================================================
\begin{titlepage}
\centering
\vspace*{1cm}

\begin{cajaTitulo}[width=\textwidth, halign=center]
  SEGURIDAD DE LA INFORMACI\'ON Y DE SISTEMAS\\[4pt]
  {\large Licenciatura en Gesti\'on de Negocios Digitales --- UCA}
\end{cajaTitulo}

\vspace{1.5cm}

{\Huge\bfseries\color{azulUCA} CUADERNILLO DE ESTUDIO}\\[8pt]
{\LARGE\color{doradoUCA} CLASE 3}\\[12pt]
{\Large\itshape ``Panorama de amenazas y\\[2pt] primeras herramientas de ciberseguridad''}\\[1.5cm]

\begin{tabular}{rl}
\textbf{Profesor:}   & Mg.\ Ing.\ Rodrigo Atilio Elgueta \\[4pt]
\textbf{Unidad:}     & Unidad 1 \\[4pt]
\end{tabular}

\vfill

\begin{tcolorbox}[colback=grisClaro, colframe=azulUCA, arc=3mm, width=0.9\textwidth]
  \centering
  {\bfseries Datos del/la estudiante}\\[10pt]
  Nombre y apellido: \underline{\hspace{5cm}}\\[10pt]
  Grupo N.\textdegree: \underline{\hspace{5cm}}\\[10pt]
  Correo: \underline{\hspace{5cm}}\\[10pt]
  Fecha: \underline{\hspace{5cm}}
\end{tcolorbox}

\vspace{1cm}
{\small\textcolor{gray}{Material de uso exclusivo para la cursada.}}
\end{titlepage}

% ============================================================
%  ÍNDICE
% ============================================================
\tableofcontents
\newpage

% ============================================================
%  SECCIÓN 1 — PANORAMA DE AMENAZAS
% ============================================================
\section{Panorama de amenazas para negocios digitales}

En la clase anterior definimos qu\'e es una \textbf{amenaza}: un evento o actor potencial que puede causar da\~no a un activo. Hoy vamos a conocer a los ``villanos'' m\'as comunes que enfrentan los e-commerce, fintech y startups en la actualidad.

% --- 1.1 Phishing ---
\subsection{Phishing y Spear Phishing}
\begin{cajaAmenaza}
El \textbf{Phishing} es el env\'io de comunicaciones masivas (correos, SMS, WhatsApp) que suplantan la identidad de una entidad confiable (un banco, Mercado Pago, Netflix) para robar credenciales o datos financieros.\\[4pt]
El \textbf{Spear Phishing} es su versi\'on dirigida: el atacante investiga a la v\'ictima en redes sociales (LinkedIn, Instagram) para crear un mensaje hiper-personalizado y mucho m\'as dif\'icil de detectar.
\end{cajaAmenaza}

\textbf{Pilar CID m\'as afectado:} Confidencialidad (robo de credenciales).

% --- 1.2 Ransomware ---
\subsection{Ransomware}
\begin{cajaAmenaza}
Un software malicioso que \textbf{cifra los archivos} del servidor o las bases de datos de la empresa y exige un rescate (generalmente en criptomonedas) para devolver la clave de descifrado. Hoy suele incluir la amenaza de publicar los datos robados si no se paga (doble extorsi\'on).
\end{cajaAmenaza}

\textbf{Pilar CID m\'as afectado:} Disponibilidad (el negocio se detiene) e Integridad/Confidencialidad (si hay fuga de datos).

% --- 1.3 Ingeniería Social ---
\subsection{Ingenier\'ia Social}
\begin{cajaAmenaza}
El arte de manipular a las personas para que revelen informaci\'on confidencial o realicen acciones que comprometan la seguridad. No ataca a la m\'aquina, \textbf{ataca al eslab\'on m\'as d\'ebil: el humano}.\\[4pt]
\emph{Ejemplo:} Llamar por tel\'efono haci\'endose pasar por ``Soporte T\'ecnico de AWS'' y pedirle a un empleado junior que le dicte el c\'odigo de autenticaci\'on de dos factores (2FA).
\end{cajaAmenaza}

% --- 1.4 Fraude / BEC ---
\subsection{Fraude y Business Email Compromise (BEC)}
\begin{cajaAmenaza}
El atacante se hace pasar por un directivo, proveedor o cliente leg\'itimo para inducir una transferencia de dinero o una acci\'on fraudulenta. En el caso BEC, se compromete el correo corporativo de un ejecutivo para enviar instrucciones falsas al \'area de finanzas.
\end{cajaAmenaza}

\textbf{Pilar CID m\'as afectado:} Confidencialidad e Integridad.

% --- 1.5 Caso de Estudio ---
\begin{cajaCaso}[title={Caso: La ca\'ida de ``FintechYa''}]
\textbf{Contexto:} Una fintech argentina de 20 empleados sufre un ataque de \emph{Spear Phishing}. El atacante env\'ia un correo al Gerente de RR.HH. con un PDF adjunto llamado ``Recibos\_Sueldo\_Actualizados.pdf.exe''.\\[4pt]
\textbf{El error:} Windows ocultaba la extensi\'on real ``.exe'' y el gerente hizo doble clic.\\[4pt]
\textbf{El impacto:} Se instal\'o un \emph{Ransomware}. La base de datos de clientes qued\'o cifrada. La empresa estuvo 5 d\'ias sin poder procesar transacciones.\\[4pt]
\textbf{Consecuencias:}
\begin{itemize}[leftmargin=1.5em]
    \item \textbf{Econ\'omico:} P\'erdida de u\$s 50.000 en ingresos no percibidos.
    \item \textbf{Reputacional:} Tendencia negativa en redes sociales por falta de respuesta.
    \item \textbf{Legal:} Denuncia ante la AAIP por no proteger los datos personales.
\end{itemize}
\end{cajaCaso}

\newpage
% ============================================================
%  SECCIÓN 2 — HERRAMIENTAS DE CIBERSEGURIDAD
% ============================================================
\section{Herramientas gratuitas de verificaci\'on}

La seguridad no es solo teor\'ia; requiere herramientas. Hoy vamos a probar tres recursos gratuitos que todo profesional de negocios digitales debe conocer.

% --- HIBP ---
\subsection{Have I Been Pwned (HIBP)}
\begin{cajaHerramienta}
\textbf{Web:} \texttt{haveibeenpwned.com}\\[4pt]
Creada por el experto Troy Hunt, esta base de datos recopila miles de millones de cuentas comprometidas en brechas de seguridad p\'ublicas. Te permite ingresar un correo electr\'onico o tel\'efono y saber si tu informaci\'on estuvo expuesta en alguna filtraci\'on hist\'orica (ej. LinkedIn, Adobe, Canva).
\end{cajaHerramienta}

% --- VirusTotal ---
\subsection{VirusTotal}
\begin{cajaHerramienta}
\textbf{Web:} \texttt{virustotal.com}\\[4pt]
Servicio propiedad de Google que analiza archivos y URLs sospechosas utilizando \textbf{m\'as de 70 motores de antivirus} simult\'aneamente (Kaspersky, Symantec, Microsoft, etc.) y servicios de sandboxing. Es ideal para verificar un enlace antes de hacer clic o un adjunto antes de abrirlo.
\end{cajaHerramienta}

% --- Gestores de Contraseñas ---
\subsection{Gestores de Contrase\~nas}
\begin{cajaHerramienta}
\textbf{Ejemplos:} Bitwarden, 1Password, KeePass.\\[4pt]
El cerebro humano no puede recordar 50 contrase\~nas distintas, complejas y \'unicas. Un gestor genera, almacena y autocompleta contrase\~nas robustas en una b\'oveda cifrada. Solo necesitas recordar \textbf{una} contrase\~na maestra.
\end{cajaHerramienta}

\begin{cajaNota}
\textbf{Regla de oro:} Nunca uses la misma contrase\~na para tu correo corporativo/personal y para un servicio de terceros (ej. un foro o un e-commerce). Si el servicio de terceros es hackeado, los atacantes probar\'an esa misma clave en tu correo (t\'ecnica de \emph{Credential Stuffing}).
\end{cajaNota}

\newpage
% ============================================================
%  SECCIÓN 3 — ACTIVIDAD PRÁCTICA (ESTACIONES)
% ============================================================
\section{Estaciones Rotativas: Manos a la obra}

En grupos, van a rotar por tres ``estaciones'' (o usar sus dispositivos m\'oviles/notebooks si el aula lo permite). En cada estaci\'on deben completar el siguiente registro.

\begin{cajaActividad}[title={Estaci\'on 1: Radiograf\'ia de Brechas (HIBP)}]
1. Ingresen a \texttt{haveibeenpwned.com}.\\
2. Busquen un correo electr\'onico antiguo o secundario (por privacidad, eviten usar su correo institucional principal en la pantalla compartida).\\
3. Respondan:
\vspace{4pt}

\renewcommand{\arraystretch}{1.5}
\begin{tabularx}{\textwidth}{|X|}
\hline
\textbf{¿En cu\'antas brechas de datos (``pwned'') apareci\'o el correo?} \\[1.5cm]
\hline
\textbf{Mencionen una de las empresas filtradas y qu\'e tipo de datos se robaron (ej. contrase\~nas, tel\'efonos, direcciones).} \\[2.5cm]
\hline
\textbf{Reflexi\'on de negocio: Si un empleado de su empresa usa ese mismo correo y contrase\~na para acceder al panel de administraci\'on del e-commerce, ¿qu\'e amenaza se materializa?} \\[2.5cm]
\hline
\end{tabularx}
\end{cajaActividad}

\vspace{8pt}

\begin{cajaActividad}[title={Estaci\'on 2: El Esc\'aner (VirusTotal)}]
1. Ingresen a \texttt{virustotal.com}.\\
2. Vayan a la pesta\~na ``URL''.\\
3. Analicen el siguiente enlace de prueba (simulado por la c\'atedra) o uno provisto por el docente:\\
\texttt{http://testsafebrowsing.appspot.com/apiv4/ANY/https/}\\
\texttt{malware.testing.google.test/testing/phishing/}\\
4. Respondan:
\vspace{4pt}

\renewcommand{\arraystretch}{1.5}
\begin{tabularx}{\textwidth}{|X|}
\hline
\textbf{¿Cu\'antos motores de seguridad marcaron la URL como maliciosa o phishing?} \\[1.5cm]
\hline
\textbf{¿Qu\'e categor\'ias o etiquetas le asignaron los antivirus (ej. Phishing, Malware, Suspicious)?} \\[2cm]
\hline
\textbf{Reflexi\'on de negocio: ¿Por qu\'e es \'util integrar una API como la de VirusTotal en el sistema de atenci\'on al cliente de una fintech cuando los usuarios reportan enlaces sospechosos?} \\[2.5cm]
\hline
\end{tabularx}
\end{cajaActividad}

\vspace{8pt}

\begin{cajaActividad}[title={Estaci\'on 3: B\'oveda de Credenciales (Gestor)}]
1. Abran un generador de contrase\~nas (puede ser el de \texttt{bitwarden.com/password-generator/} o uno local).\\
2. Config\'urenlo para generar una clave de: \textbf{16 caracteres, con may\'usculas, min\'usculas, n\'umeros y s\'imbolos}.\\
3. Respondan:
\vspace{4pt}

\renewcommand{\arraystretch}{1.5}
\begin{tabularx}{\textwidth}{|X|}
\hline
\textbf{Copien la contrase\~na generada (o una parte de ella):} \\[1.5cm]
\hline
\textbf{¿Cu\'anto tiempo estiman que tardar\'ia un ataque de fuerza bruta en adivinar esta clave comparado con la clave ``Empresa123''?} \\[2cm]
\hline
\textbf{Reflexi\'on de negocio: ¿Qu\'e pol\'itica de seguridad implantar\'ian en su startup respecto al uso de gestores de contrase\~nas para el equipo?} \\[2.5cm]
\hline
\end{tabularx}
\end{cajaActividad}

\newpage
% ============================================================
%  SECCIÓN 4 — TP INTEGRADOR
% ============================================================
\section{Conformaci\'on de Grupos y TP Integrador}

A partir de hoy, trabajar\'an en grupos de 4 a 5 personas. Este grupo los acompa\~nar\'a hasta el final de la cursada para desarrollar el \textbf{Trabajo Pr\'actico Integrador: ``Programa de Seguridad para mi Negocio Digital''}.

\subsection{¿En qu\'e consiste el TP Integrador?}
Elegir\'an (o se les asignar\'a) un negocio digital (e-commerce, fintech, app, marketplace) y construir\'an su dossier de seguridad con 6 entregables:
\begin{enumerate}[leftmargin=2em]
    \item \textbf{Mapa de activos} y an\'alisis CID (Unidad 1).
    \item \textbf{Matriz de riesgos} y BIA simplificado (Unidad 2).
    \item Evaluaci\'on de seguridad en la \textbf{nube}.
    \item Cumplimiento normativo y \textbf{pol\'itica de privacidad} (Unidad 3).
    \item \textbf{Pol\'itica de Seguridad Aceptable (PSA)} (Unidad 4).
    \item \textbf{Plan de respuesta a incidentes} (verificado con IA).
\end{enumerate}
Adem\'as, un \textbf{Anexo ``Uso de IA''} documentando qu\'e partes se apoyaron en IA generativa y c\'omo se verific\'o.

\begin{cajaActividad}[title={Registro del Grupo}]
Completen los datos de su equipo y la idea de negocio que intentar\'an desarrollar (sujeta a aprobaci\'on del docente).

\vspace{6pt}
\renewcommand{\arraystretch}{1.6}
\begin{tabularx}{\textwidth}{|>{\bfseries}p{4.5cm}|X|}
\hline
\textbf{Nombre del Grupo / Empresa:} & \hrulefill \\[6pt]
\hline
\textbf{Integrante 1:} & \hrulefill \\[6pt]
\hline
\textbf{Integrante 2:} & \hrulefill \\[6pt]
\hline
\textbf{Integrante 3:} & \hrulefill \\[6pt]
\hline
\textbf{Integrante 4:} & \hrulefill \\[6pt]
\hline
\textbf{Integrante 5 (si aplica):} & \hrulefill \\[6pt]
\hline
\textbf{Idea de Negocio Digital:} & \\
\multicolumn{2}{|p{\dimexpr\textwidth-2\tabcolsep-2\arrayrulewidth}|}{\vspace{0.3cm} \hrulefill \vspace{0.3cm} \hrulefill \vspace{0.3cm} \hrulefill \vspace{0.3cm}} \\[6pt]
\hline
\end{tabularx}
\end{cajaActividad}

\newpage
% ============================================================
%  SECCIÓN 5 — GLOSARIO Y CIERRE
% ============================================================
\section{Glosario de la Clase 3}

\renewcommand{\arraystretch}{1.4}
\begin{tabularx}{\textwidth}{>{\bfseries}p{4.5cm}X}
\toprule
\textbf{T\'ermino} & \textbf{Definici\'on en una l\'inea} \\
\midrule
Phishing & Suplantaci\'on de identidad para robar credenciales. \\
Spear Phishing & Phishing dirigido e hiper-personalizado. \\
Ransomware & Malware que secuestra datos cifr\'andolos y pide rescate. \\
Ingenier\'ia Social & Manipulaci\'on psicol\'ogica del factor humano. \\
BEC (Business Email Compromise) & Fraude mediante suplantaci\'on de correo corporativo. \\
Credential Stuffing & Ataque automatizado usando contrase\~nas filtradas en otros sitios. \\
Breach (Brecha) & Incidente donde la informaci\'on confidencial es expuesta. \\
Sandbox & Entorno aislado para ejecutar archivos sospechosos sin riesgo. \\
2FA / MFA & Autenticaci\'on de dos o m\'ultiples factores (algo que s\'e + algo que tengo). \\
\bottomrule
\end{tabularx}

\vspace{1cm}

\section{Espacio para notas y hallazgos}

\begin{tcolorbox}[colback=grisClaro, colframe=gray!50, arc=2mm, height=6cm]
\textcolor{gray}{\itshape Anot\'a aqu\'i los hallazgos m\'as interesantes de las herramientas, ideas para tu TP Integrador o dudas para la pr\'oxima clase sobre criptograf\'ia.}
\end{tcolorbox}

\vfill
\begin{center}
\small\textcolor{gray}{%
Cuadernillo elaborado para la c\'atedra de Seguridad de la Informaci\'on y de
Sistemas --- Lic.\ en Gesti\'on de Negocios Digitales, UCA.\\
Prohibida su distribuci\'on fuera del \'ambito de la cursada sin autorizaci\'on del docente.}
\end{center}

\end{document}